\section{Introduction}
Diffusion models have achieved remarkable success in image \cite{sdxl, sana} and video \cite{stablevideo, hunyuanvideo} generation, and have recently been extended to text generation. Unlike autoregressive (AR) models \cite{gpt4, qwen2.5} that generate tokens sequentially, diffusion large language models (dLLMs) \cite{llada, dream} adopt full attention over all positions and refine an entire sequence through multiple denoising iterations, achieving surprisingly strong performance on text generation tasks. dLLMs offer potential solutions to longstanding limitations of AR models, including the inability to decode in parallel and the reversal curse \cite{reversal}. These characteristics have drawn growing research interest in further advancing dLLMs \cite{gemini_diffusion, seeddiffusion, mercury, llada2.0}. 
% More recently, efforts \cite{llada2.0, wedlm} to scale and improve dLLMs have further advanced this emerging paradigm.

Despite these efforts, dLLMs still exhibit performance gaps in practical deployments compared to state-of-the-art AR models \cite{fastdllm, eagle3}. These gaps are largely attributed to two main factors: KV-Cache management \cite{fastdllm, dkvcache, dllmcache} and nonindependent position predictions \cite{parallelcurse, fastdllm}. First, dLLMs employ bidirectional attention, which fundamentally contradicts the causal assumption underlying standard KV-Cache mechanisms. Second, the marginal distributions at each position produced by dLLMs often violate the independence assumption underlying parallel decoding.

This work aims to improve the efficiency of parallel decoding in dLLMs, with a particular focus on \textit{nonindependent position predictions}. Most existing parallel decoding strategies \cite{fastdllm, ebsampler} select masked positions using heuristics and relatively coarse-grained criteria, such as confidence and entropy, to ensure that the selected positions behave approximately independently. However, overly conservative selection criteria can substantially limit the achievable parallelism, leaving much of the potential efficiency of parallel decoding underexploited. To address this problem, a natural alternative is to account for positional relationships more directly. Since the main difficulty stems from position coupling, improving parallel decoding requires approximating positional dependencies during inference rather than evaluating each position in isolation. Attention maps \cite{spargeattn, sageattention} provide an approximate yet cheap signal of token interactions that is readily available from each forward pass, making them a practical proxy for dependency estimation. From this perspective, two observations are particularly relevant: (i) we verify that dLLMs can exhibit abnormal attention concentration patterns (akin to attention sinks \cite{streamllm, attentionsinksdiffusionlanguage}) that are largely semantically irrelevant, which mislead the attention-based dependency estimates; and (ii) positions that are strongly dependent to previously unmasked high-confidence tokens can remain highly consistent with the final output even when their confidence is relatively low. Moreover, 
% high-confidence positions are approximately independent in practice, whereas other positions can be strongly coupled; consequently, 
avoiding simultaneous unmasking of strongly coupled low-confidence positions can substantially reduce errors induced by parallel decoding under marginal probabilities. These observations provide a new insight: \textit{dependency-aware inference rules can expand safe parallelism beyond those conservative threshold methods}.

Motivated by this, we propose \dawn{}, a training-free Dependency-AWare fast inference method for diffusioN LLMs. It improves parallel decoding by explicitly accounting for position dependency. 
\dawn{} consists of three cooperating components: \textit{Dependency Graph Construction}, \textit{Anchor-Guided Decoding}, and \textit{Conflict-Based Scheduling}. At each denoising iteration, a dependency graph is constructed from processed attention maps via thresholding, which captures salient token coupling relations. Based on this graph, the subsequent two components select two disjoint sets of positions for parallel updates, $\mathcal{U}_{\text{anchor}}$ and $\mathcal{U}_{\text{conflict}}$. \textit{Anchor-Guided Decoding} first selects approximately independent high-confidence positions, and then treats previously unmasked high-confidence tokens as anchors and relaxes the confidence threshold for strongly anchor-coupled masked positions. \textit{Conflict-Based Scheduling} identifies conflicts in the dependency graph and greedily constructs a large non-conflicting update set from the remaining candidates that satisfy a lower confidence threshold. Positions in $\mathcal{U}^{(t)}_{\text{anchor}} \cup~ \mathcal{U}^{(t)}_{\text{conflict}}$ are then unmasked in parallel. 

Compared with prior approaches that rely on strict positional criteria, \dawn{} relaxes selection constraints by incorporating dependency information, thereby enabling additional unmasking that would otherwise be filtered out while preserving generation quality. Meanwhile, it substantially mitigates failures caused by nonindependent position predictions during parallel decoding.  Extensive experiments validate that \dawn{} improves the quality--speed trade-off of dLLM inference across multiple models and settings.

The main contributions are summarized as follows:
\begin{itemize}[leftmargin=*, itemsep=2pt, parsep=2pt, topsep=0pt, partopsep=0pt]

    % \item \todo[inline]{new strategies, observations}\dawn{} uses attention maps to approximately characterize dependencies among positions during inference, while correcting spurious signals caused by attention sinks. Building on these dependency cues, we further identify characteristic unmasking behaviors of coupled tokens under different conditions and design corresponding decoding rules that yield strong empirical performance.

    % \item We directly extract inter-position dependencies during inference and leverage them to alleviate a core difficulty in dLLM parallel decoding: non-independent positional predictions. Motivated by this perspective, we verify that attention sinks can bias the attention-based dependency proxies, and observe that previously unmasked high-confidence tokens can reliably induce unmasking at dependent positions. These findings achieve dependency-aware decoding rules that achieve strong empirical performance.

    \item We show that token dependencies can be estimated from attention maps during inference, enabling more aggressive parallel decoding. We identify two key findings: (1) attention sinks—positions that absorb disproportionate attention regardless of semantics—distort dependency estimates and must be filtered; (2) once a high-confidence token is committed, positions that depend on it become reliably predictable, even at lower confidence.
    
    % \item We propose \textbf{\dawn{}}, a training-free, dependency-aware method for fast inference in diffusion large language models. \dawn{} comprises three components: \textit{Dependency Graph Construction}, \textit{Anchor-Guided Decoding}, and \textit{Conflict-Based Scheduling}, and mitigates the nonindependent issue in dLLM parallel decoding, thereby enabling more efficient inference.

    \item We propose \textbf{\dawn{}}, a training-free, dependency-aware method for fast inference of diffusion LLMs. It 
    %comprises three components: \textit{Dependency Graph Construction}, \textit{Anchor-Guided Decoding}, and \textit{Conflict-Based Scheduling}. \dawn{} 
    uses the estimated dependencies in two ways: relaxing confidence thresholds for positions anchored by committed high-confidence tokens, and preventing strongly coupled low-confidence positions from being unmasked together, thereby enabling more efficient inference.
    
    \item Extensive experiments across multiple models, datasets, and representative baselines demonstrate the effectiveness of \dawn{}, achieving \textbf{1.80 - 8.06$\times$} speedup over the baseline. Ablation studies further validate the contributions of each component and the impact of key hyperparameters. 
\end{itemize}
% 1: We show that token dependencies can be estimated from attention maps during inference, enabling more aggressive parallel decoding. We identify two key findings: (1) attention sinks—positions that absorb disproportionate attention regardless of semantics—distort dependency estimates and must be filtered; (2) once a high-confidence token is committed, positions that depend on it become reliably predictable, even at lower confidence.

% 2: DAWN uses the estimated dependencies in two ways: it relaxes confidence thresholds for positions anchored by committed high-confidence tokens, and it prevents strongly coupled low-confidence positions from being unmasked together.

